\VerdierPhysicalPageAnchor{0057}%
\VerdierSemanticAnchor{ch1-1-5-7-continuation}
fait apparaître cinq diagrammes verticaux qui sont respectivement les
diagrammes :

\VerdierSemanticAnchor{ch1-1-5-7-vertical}
\[
  T'(\sigma')\Delta(\sigma,\tau),\quad
  T'(\sigma)\Delta(\sigma',\tau),\quad
  \Delta(\sigma,\tau)T(\sigma'),\quad
  \Delta(\sigma',\tau)T(\sigma),\quad
  \Delta(\sigma+\sigma',\tau).
\]

Il fait apparaître de même quatre diagrammes horizontaux qui sont
respectivement les diagrammes :

\VerdierSemanticAnchor{ch1-1-5-7-horizontal}
\[
  D(\sigma,\sigma')T(\tau),\quad
  D(\sigma',\sigma)T(\tau),\quad
  T'(\tau)D(\sigma,\sigma'),\quad
  T'(\tau)D(\sigma',\sigma).
\]

\VerdierSemanticAnchor{ch1-1-5-7-conclusion}
Par hypothèse de récurrence, les diagrammes $D(\sigma,\sigma')$ et
$D(\sigma',\sigma)$ sont respectivement $\alpha(\sigma,\sigma')$ et
$\alpha(\sigma',\sigma)$-commutatifs. On en déduit la première partie
de l'assertion (1.5.4.3) en envisageant les trois cas possibles. La
deuxième partie de l'assertion (1.5.4.3) s'obtient en interchangeant les
rôles des $\sigma$ et des $\tau$. Dans le cas général, le diagramme
qu'il faut introduire est plus compliqué. Mais, en utilisant les
identités (1.2.1.3), l'argument se ramène à celui présenté ci-dessus.

\VerdierSemanticAnchor{ch1-1-5-8}
\noindent\textbf{Remarque 1.5.8.}\quad
La proposition (1.5.2) s'applique en particulier au cas où le cocycle
$\alpha$ est trivial. On vérifie alors que la proposition et sa
démonstration s'étendent \emph{verbatim} au cas où la catégorie
$\mathcal C'$ n'est plus nécessairement additive.

\VerdierSemanticAnchor{ch1-1-5-9}
\noindent\textbf{1.5.9.}\quad
Soient :

\VerdierSemanticAnchor{ch1-1-5-9-functors}
\[
  (F_1,p_1),(F_2,p_2)\colon
  (\mathcal C,\mathbb Z^n,T,c)
  \longrightarrow
  (\mathcal C',\mathbb Z^n,T',c')
\]
deux $\mathbb Z^n$-foncteurs tordus respectivement par deux cocycles
normalisés $\alpha_1,\alpha_2\in
Z^2(\mathbb Z^n,\mathbb Z/2\mathbb Z)$,
$\beta\colon\mathbb Z^n\to\mathbb Z/2\mathbb Z$ une application
telle que :

\VerdierSemanticAnchor{ch1-1-5-9-eq1}
\begin{equation}
  \alpha_2(\sigma_1,\sigma_2)\alpha_1^{-1}(\sigma_1,\sigma_2)
  =\beta(\sigma_2)\beta(\sigma_1\sigma_2)^{-1}\beta(\sigma_1),
  \qquad \sigma_1,\sigma_2\in\mathbb Z^n,
  \tag{1.5.9.1}
\end{equation}
et $u\colon(F_1,p_1)\to(F_2,p_2)$ un $\beta$-morphisme de foncteurs
(1.4.6). Soit $e_1,\ldots,e_n$ la base canonique de $\mathbb Z^n$.
Pour tout $i\in[1,n]$ le diagramme ci-après est
$\beta(e_i)$-commutatif :

\VerdierSemanticAnchor{ch1-1-5-9-eq2}
\begin{equation}
  \resizebox{0.78\textwidth}{!}{%
    $\begin{tikzcd}[ampersand replacement=\&,column sep=6em,row sep=4em]
      F_1T(e_i)
        \arrow[r,"{p_1(e_i)}"]
        \arrow[d,"{u*T(e_i)}"']
        \& T'(e_i)F_1
          \arrow[d,"{T'(e_i)*u}"] \\
      F_2T(e_i)
        \arrow[r,"{p_2(e_i)}"]
        \& T'(e_i)F_2
    \end{tikzcd}$%
  }\;.
  \tag{1.5.9.2}
\end{equation}
